\uuid{y3f2}
\titre{Famille libre/Génératrice/Base}
\auteur{hanani ; mzari}
\organisation{exo7}
\licence{CC BY-NC-SA 4.0}
\datecreate{2020-01-20}
\niveau{L1}
\module{Algèbre}
\chapitre{Espace vectoriel}
\sousChapitre{Base}
\difficulte{1}
\theme{Base et dimension, Famille libre, Génératrice, Base}
\exo7qcmid{lille-2.espaces-vectoriels.027}

\contenu{
\begin{question}
\qclassification{106.04}
\qtags{motcle=Famille libre/Génératrice/Base}
Dans $\Rr^3$, on considère les vecteurs $u_1=(1,1,1), u_2=(0,1,1)$ et $ u_3=(-1,1,0)$.
Quelles sont les assertions vraies ?
\begin{answers}
\good{$\{u_1,u_2,u_3\}$ est une famille libre}.
\good{$\{u_1,u_2,u_3\}$ est une famille génératrice de $\Rr^3$}.
\bad{$u_2$ est une combinaison linéaire de $u_1$ et $u_3$}.
\bad{$\{u_1,u_2,u_3\}$ n'est pas une base de $\Rr^3$}.
\end{answers}
\begin{explanations} On vérifie que $\{u_1,u_2,u_3\}$ est une famille libre. Comme cette famille contient $3$  vecteurs
linéairement indépendants de $\Rr^3$ et la dimension de $\Rr^3$ est $3$, elle est génératrice de   $\Rr^3$ et donc c'est une base de $\Rr^3$.
\end{explanations}
\end{question}
}
